\documentclass[11pt]{article}

% Packages

\usepackage[margin=1in]{geometry}

% Figures
\usepackage{graphicx}
\usepackage{float}

% Tables
\usepackage{booktabs}
\usepackage{array}

% Mathematics
\usepackage{amsmath}
\usepackage{amssymb}

% Units and number formatting
\usepackage{siunitx}

% Better links
\usepackage[hidelinks]{hyperref}

% Better spacing
\usepackage{setspace}

% Optional code blocks
\usepackage{listings}

% General formatting

\setstretch{1.05}

\sisetup{
    group-separator={,},
    group-minimum-digits=4
}

% Reusable commands

% Project name
\newcommand{\project}{Stratax}

% Benchmark suite version
\newcommand{\benchmarkversion}{V2}

% Number of warmups
\newcommand{\warmups}{10}

% Number of timed samples per trial
\newcommand{\samples}{50}

% Number of trials
\newcommand{\trials}{5}

% Document information

\title{
    \textbf{\project{} \benchmarkversion{} Benchmark Report}\\
    \large Cross-Platform and Compiler Performance Analysis
}

\author{Noah Ryan}

\date{\today}


\begin{document}

\maketitle

% Abstract

\begin{abstract}

This report evaluates the performance characteristics of
\project{}, a scientific array library implemented in C++.
Benchmark Suite \benchmarkversion{} measures representative
array operations across multiple hardware, operating-system,
and compiler configurations.

Each benchmark configuration uses repeated warmup executions,
timed samples, and independent trials to measure both runtime
and measurement variability.

The purpose of this report is to document the benchmark
methodology, present the resulting measurements, identify
performance trends, and establish a reproducible baseline for
future optimization and regression testing.

\end{abstract}


% Introduction / Objective
\section{Objective}

The primary objective of Benchmark Suite \benchmarkversion{}
is to characterize the runtime performance of \project{} under
different execution environments.

The experiment considers several questions:

\begin{itemize}
    \item How does runtime scale as the number of array elements increases?
    \item How much does compiler choice affect performance on the same processor?
    \item How does the operating environment affect performance?
    \item Are performance differences consistent across different operations?
    \item How reproducible are repeated benchmark measurements?
\end{itemize}

The results are intended to provide a baseline against which
future changes to \project{} can be compared.

% Methodology

\section{Methodology}

\subsection{Benchmark Procedure}

Each benchmark case follows the same general procedure.

First, an untimed correctness validation is performed.
The benchmark operation is then executed \warmups{} times as
a warmup before measurements begin.

Each trial contains \samples{} timed executions of the operation.
A total of \trials{} independent trials are performed for each
platform and compiler configuration.

Timing is performed using:

\begin{center}
\texttt{std::chrono::steady\_clock}
\end{center}

The timed interval contains the benchmarked operation and any
result allocation performed by that operation.

Input construction, correctness validation, checksum computation,
result destruction, and console output are excluded from the
timed interval.


\subsection{Problem Sizes}

The default benchmark suite uses three square array sizes.

\begin{table}[H]
    \centering
    \begin{tabular}{rr}
        \toprule
        \textbf{Dimensions} & \textbf{Elements} \\
        \midrule
        $64 \times 64$     & \num{4096} \\
        $256 \times 256$   & \num{65536} \\
        $1024 \times 1024$ & \num{1048576} \\
        \bottomrule
    \end{tabular}
    \caption{Array sizes used by Benchmark Suite \benchmarkversion{}.}
    \label{tab:sizes}
\end{table}


\subsection{Benchmark Configuration}

The default benchmark parameters are summarized in
Table~\ref{tab:configuration}.

\begin{table}[H]
    \centering
    \begin{tabular}{lr}
        \toprule
        \textbf{Parameter} & \textbf{Value} \\
        \midrule
        Benchmark suite version & \benchmarkversion{} \\
        Warmup executions       & \warmups{} \\
        Timed samples per trial & \samples{} \\
        Independent trials      & \trials{} \\
        Primary timing statistic & Median \\
        Clock & \texttt{std::chrono::steady\_clock} \\
        \bottomrule
    \end{tabular}
    \caption{Benchmark methodology parameters.}
    \label{tab:configuration}
\end{table}

% Experimental systems

\section{Experimental Platforms}

Benchmark Suite \benchmarkversion{} was executed on multiple
hardware and software configurations.

\begin{table}[H]
    \centering
    \begin{tabular}{llll}
        \toprule
        \textbf{Processor} &
        \textbf{OS} &
        \textbf{Compiler} &
        \textbf{Configuration} \\
        \midrule

        Intel Core Ultra 9 288V &
        Debian &
        GCC &
        Laptop \\

        Intel Core Ultra 9 288V &
        Debian &
        Clang &
        Laptop \\

        Ryzen 9 9950X &
        Ubuntu &
        GCC &
        Desktop \\

        Ryzen 9 9950X &
        Ubuntu &
        Clang &
        Desktop \\

        Ryzen 9 9950X &
        Windows &
        Clang &
        Desktop \\

        Ryzen 9 9950X &
        Windows &
        MSVC &
        Desktop \\

        \bottomrule
    \end{tabular}

    \caption{
        Hardware, operating-system, and compiler configurations.
    }

    \label{tab:platforms}
\end{table}

Exact compiler versions and additional environment metadata are
stored alongside the benchmark results in the corresponding JSON
files.

% Benchmark workloads

\section{Benchmark Workloads}

Benchmark Suite \benchmarkversion{} contains several operations
intended to exercise different parts of the \project{} array
implementation.


\subsection{Iteration}

The iteration benchmarks measure traversal over owning arrays
and array views.

\begin{description}

    \item[\texttt{owning\_iteration}]
    Sequential iteration over an owning array.

    \item[\texttt{contiguous\_view\_iteration}]
    Iteration over a contiguous array view.

    \item[\texttt{strided\_view\_iteration}]
    Iteration over a non-contiguous or strided array view.

\end{description}


\subsection{Arithmetic and Broadcasting}

These benchmarks measure elementwise operations and broadcasting.

\begin{description}

    \item[\texttt{owning\_scalar\_add}]
    Adds a scalar value to each element of an owning array.

    \item[\texttt{owning\_row\_broadcast\_add}]
    Adds a row vector to an owning matrix using broadcasting.

\end{description}


\subsection{Reductions}

Reduction benchmarks measure operations that combine many array
elements into smaller results.

\begin{description}

    \item[\texttt{owning\_axis0\_sum}]
    Sum reduction along axis 0.

    \item[\texttt{owning\_axis1\_sum}]
    Sum reduction along axis 1.

    \item[\texttt{strided\_view\_axis1\_sum}]
    Axis-1 reduction over a strided view.

\end{description}


\subsection{Materialization and Copying}

These benchmarks measure operations that produce new owning
storage from existing arrays or views.

\begin{description}

    \item[\texttt{contiguous\_view\_materialize}]
    Materializes a contiguous view into owning storage.

    \item[\texttt{strided\_view\_materialize}]
    Materializes a strided view into owning storage.

    \item[\texttt{owning\_to\_tensor\_copy}]
    Copies an owning array into a tensor representation.

\end{description}

% Statistics

\section{Statistical Analysis}

Performance measurements are affected by operating-system
scheduling, interrupts, background activity, cache state, and
other sources of variability. Multiple statistics are therefore
used to describe the measurements.


\subsection{Median}

The median is used as the primary timing statistic because it is
less sensitive to isolated timing spikes than the arithmetic mean.

For an ordered set of measurements

\[
x_1, x_2, \ldots, x_n,
\]

the median represents the central measurement.


\subsection{Arithmetic Mean}

The arithmetic mean is calculated as

\begin{equation}
    \bar{x}
    =
    \frac{1}{n}
    \sum_{i=1}^{n} x_i.
    \label{eq:mean}
\end{equation}


\subsection{Sample Standard Deviation}

The sample standard deviation is

\begin{equation}
    s
    =
    \sqrt{
        \frac{
            \sum_{i=1}^{n}
            \left(x_i - \bar{x}\right)^2
        }{
            n - 1
        }
    }.
    \label{eq:stdev}
\end{equation}

This statistic measures the spread of repeated timing
measurements around their mean.


\subsection{Median Absolute Deviation}

Median absolute deviation, or MAD, is calculated as

\begin{equation}
    \mathrm{MAD}
    =
    \mathrm{median}
    \left(
        \left|
            x_i - \mathrm{median}(x)
        \right|
    \right).
    \label{eq:mad}
\end{equation}

MAD provides a measure of spread that is less sensitive to
outliers than standard deviation.


\subsection{Coefficient of Variation}

The coefficient of variation is

\begin{equation}
    \mathrm{CV}
    =
    \frac{s}{\bar{x}}
    \times 100\%.
    \label{eq:cv}
\end{equation}

CV expresses variability relative to the mean runtime, making it
easier to compare noise between benchmarks with different
absolute runtimes.


\subsection{Across-Trial Statistics}

Each benchmark configuration is executed for \trials{}
independent trials.

The plotted runtime is the median of the individual trial
medians:

\begin{equation}
    T
    =
    \mathrm{median}
    \left(
        m_1,
        m_2,
        \ldots,
        m_{\trials}
    \right),
    \label{eq:trialmedian}
\end{equation}

where $m_i$ is the median runtime from trial $i$.

Error bars represent the sample standard deviation of the
trial medians.


\subsection{Relative Performance}

When comparing two configurations, relative speedup may be
calculated as

\begin{equation}
    S
    =
    \frac{
        T_{\mathrm{baseline}}
    }{
        T_{\mathrm{comparison}}
    }.
    \label{eq:speedup}
\end{equation}

A value of

\[
S > 1
\]

indicates that the comparison configuration completed the
operation in less time than the baseline.

% Results

\section{Results}

This section presents selected benchmark results.
The complete set of generated plots is available in the
\texttt{benchmarks/plots/} directory.


\subsection{Iteration Performance}

Figure~\ref{fig:contiguous-iteration} shows an example comparison
of contiguous view iteration across the tested configurations.

\begin{figure}[H]
    \centering

    \includegraphics[
        width=0.90\textwidth
    ]{
        ../../plots/v2/contiguous_view_iteration.png
    }

    \caption{
        Contiguous view iteration performance.
        Each point represents the median of the trial medians.
        Error bars represent the standard deviation of trial
        medians.
    }

    \label{fig:contiguous-iteration}
\end{figure}

% Example analysis text:
%
% Describe what you observe here.
%
% Avoid simply saying:
% "Configuration A is better."
%
% Instead explain:
% - how performance changes with size
% - whether configurations remain separated
% - whether differences grow or shrink
% - whether error bars indicate stable measurements


\subsection{Arithmetic Performance}

% Example figure:
%
% \begin{figure}[H]
%     \centering
%     \includegraphics[
%         width=0.90\textwidth
%     ]{
%         ../../plots/v2/owning_scalar_add.png
%     }
%     \caption{Owning scalar addition performance.}
%     \label{fig:scalar-add}
% \end{figure}


\subsection{Broadcasting Performance}

% Discuss owning_row_broadcast_add.


\subsection{Reduction Performance}

% Discuss axis0 and axis1 reduction behavior.


\subsection{Materialization and Copy Performance}

% Discuss materialization and tensor-copy behavior.

% Example statistical table

\subsection{Selected Numerical Results}

Table~\ref{tab:example-results} demonstrates one way to present
a small subset of exact measurements.

\begin{table}[H]
    \centering

    \begin{tabular}{lrrr}
        \toprule
        \textbf{Configuration} &
        \textbf{Elements} &
        \textbf{Median (\si{\micro\second})} &
        \textbf{Trial CV (\%)} \\
        \midrule

        Ubuntu GCC &
        \num{1048576} &
        0.0 &
        0.0 \\

        Ubuntu Clang &
        \num{1048576} &
        0.0 &
        0.0 \\

        Windows Clang &
        \num{1048576} &
        0.0 &
        0.0 \\

        Windows MSVC &
        \num{1048576} &
        0.0 &
        0.0 \\

        \bottomrule
    \end{tabular}

    \caption{
        Example format for reporting selected benchmark statistics.
        Replace the placeholder values with measurements from the
        generated summary data.
    }

    \label{tab:example-results}
\end{table}

% Discussion

\section{Discussion}

The results should be interpreted by considering both runtime
and measurement variability.

This section should focus on patterns rather than simply
repeating values from the tables.

Useful questions include:

\begin{itemize}

    \item Does runtime scale proportionally with the number of elements?

    \item Do compiler differences remain consistent across problem sizes?

    \item Do different operations respond differently to compiler choice?

    \item Are there benchmarks where one configuration behaves
          unexpectedly?

    \item Are observed differences much larger than the measured
          trial-to-trial variation?

\end{itemize}

Large or unexpected differences should be treated as observations
that may require additional investigation rather than immediately
being assigned a cause.

% Limitations

\section{Limitations}

Benchmark Suite \benchmarkversion{} provides controlled and
repeatable measurements, but several limitations remain.

\begin{itemize}

    \item The laptop and desktop use different processor architectures,
          memory systems, thermal limits, and power-management behavior.

    \item Comparisons between Windows and Linux also involve differences
          in runtime libraries, standard-library implementations, ABI,
          and system behavior.

    \item CPU affinity is not explicitly controlled.

    \item CPU frequency and thermal state are not directly fixed.

    \item Background operating-system activity may affect timing.

    \item Extremely short operations can be difficult to measure
          precisely when timed individually.

    \item Benchmark Suite \benchmarkversion{} measures \project{}
          internally and does not compare against external numerical
          libraries such as NumPy.

\end{itemize}

% Future work

\section{Future Work}

Future benchmark studies may investigate:

\begin{itemize}

    \item multithreaded execution,
    \item SIMD/vectorization,
    \item larger problem sizes,
    \item cache and memory-bandwidth behavior,
    \item additional processors,
    \item release-to-release performance regressions,
    \item comparisons with NumPy and other numerical libraries.
\end{itemize}

% Conclusion

\section{Conclusion}

Benchmark Suite \benchmarkversion{} establishes a reproducible
performance baseline for \project{} across multiple compiler,
operating-system, and processor configurations.

The results provide both performance measurements and estimates
of measurement variability. These measurements can be used to
guide optimization work and to identify future performance
regressions.

% Add 1--2 paragraphs describing the most important findings
% from this specific experiment.

% Appendix

\appendix


\section{Benchmark Parameters}

\begin{table}[H]
    \centering

    \begin{tabular}{lr}
        \toprule
        \textbf{Parameter} & \textbf{Value} \\
        \midrule
        Benchmark suite & \benchmarkversion{} \\
        Warmups & \warmups{} \\
        Samples per trial & \samples{} \\
        Trials & \trials{} \\
        Small size & $64 \times 64$ \\
        Medium size & $256 \times 256$ \\
        Large size & $1024 \times 1024$ \\
        \bottomrule
    \end{tabular}

    \caption{Benchmark Suite \benchmarkversion{} parameters.}
\end{table}


\section{Reproduction}

A benchmark trial can be executed using the benchmark driver.

\begin{lstlisting}[language=bash]
python scripts/benchmark.py \
    --trial 1 \
    --compiler g++ \
    --output results.json
\end{lstlisting}

Statistical analysis can be regenerated using:

\begin{lstlisting}[language=bash]
python benchmarks/analyze.py
\end{lstlisting}

Plots can be regenerated using:

\begin{lstlisting}[language=bash]
python benchmarks/plot.py
\end{lstlisting}


\section{Repository Data}

The complete machine-readable benchmark results are stored in:

\begin{center}
\texttt{benchmarks/results/v2/}
\end{center}

Generated figures are stored in:

\begin{center}
\texttt{benchmarks/plots/v2/}
\end{center}

Statistical summaries are stored in:

\begin{center}
\texttt{benchmarks/analysis/v2/}
\end{center}


\end{document}